\documentclass[10pt,a4paper]{amsart}
\usepackage[margin=19mm]{geometry}
\usepackage{amsmath,amssymb,mathtools}
\usepackage[scr=rsfs,cal=euler]{mathalfa}
\usepackage{xfrac}
\usepackage[unicode,hidelinks,bookmarks=false]{hyperref}
\newcommand{\ie}{i.e.\ }
\newcommand{\cf}{cf.\ }
\newcommand{\resp}{respectively\ }
\newcommand{\Subsection}{\subsection}
\providecommand{\nobreakdash}{\nobreak\hbox{-}}
\setlength{\emergencystretch}{2em}
\multlinegap0pt
\usepackage{amssymb}
\usepackage{tikz}
\usepackage{mathtools}
\usepackage{algorithm}\usepackage{algpseudocode}
\usepackage{float}
\usepackage[shortlabels]{enumitem}
\usepackage{xcolor}
\usepackage[normalem]{ulem}
\newtheorem{lemma}{Lemma}
\newtheorem{theorem}{Theorem}
\newcommand{\R}{{\mathbb R}}
\newcommand{\zer}{\operatorname{zer}}
\title{Inexact Proximal Point and Tseng Algorithms with Nonsummable Errors to Solve Monotone Inclusions}
\author{Ba Khiet Le, Boris S. Mordukhovich, Michel A. Thera}
\date{Source: arXiv:2606.01536v2 (CC BY 4.0)}
\begin{document}
\maketitle

\noindent\emph{Source and licence.} Inexact Proximal Point and Tseng Algorithms with Nonsummable Errors to Solve Monotone Inclusions. Version arXiv:2606.01536v2 (CC BY 4.0), distributed by arXiv under the Creative Commons Attribution 4.0 International licence, http://creativecommons.org/licenses/by/4.0/, as recorded in the arXiv licence metadata for this exact version and shown on the abstract page https://arxiv.org/abs/2606.01536v2. Full licence notice: \texttt{licenses/arxiv-2606.01536-CC-BY-4.0.txt}.\par\smallskip

\noindent\textit{Editorial scope.} Counted unit: the article's theorem thm:ita (source0.tex lines 739--767). The article's complete original proof of it is delivered with it (source0.tex lines 768--827, one \texttt{proof} environment of 60 lines). No mathematics is written, repaired or re-derived: the statement and the proof are the article's own lines, in the article's own notation, and no retained line is substituted or re-flowed.  Retained uncounted as the source-local context the proof needs: lines 651--655; lines 686--691.  Nothing was omitted from any retained span, so the package carries no omission record.  No retained line contains a citation, so the wrapper carries no bibliography and no bibliography entry.  No retained line contains a figure environment, an image-inclusion command or drawing code, so no image is delivered and the package declares no asset dependency.  Editorial formatting only, in addition: an AMS \texttt{amsart} preamble that restates the article's own declarations and macros used by the retained lines, namely the environments \texttt{lemma}, \texttt{theorem} on separate counters, together with the standard packages those lines need.  The retained environments print in this document as Lemma 1, Theorem 1 in order of appearance, whereas the article prints the same items with the numbering of its own full document.  The complete native source of the article as distributed in the arXiv e-print for this exact version is bundled unchanged as \texttt{sources/201837/source0.tex}.
\begin{equation}\label{tseng1}
0\in Ax+Bx+\epsilon x
=
\widetilde{A}x+Bx,
\end{equation}

\begin{lemma}\label{ab}
Let $x^*$ be the unique fixed point of $T$. Choose $\gamma>0$ such that $\gamma^2 L^2\le 1-2\gamma\varepsilon $. Then the operator $\Vert T(\cdot)-x^*\Vert: \mathcal{H}\to \R_+$ is $\kappa$-Lipschitz continuous with the modulus
$$
\kappa=\sqrt{1-\gamma\epsilon}<1.
$$
\end{lemma}

\begin{theorem}\label{thm:ita} Let $A:\mathcal{H}\rightrightarrows \mathcal{H}$ be an arbitrary maximally monotone operator, and let the operator $B:\mathcal{H}\to\mathcal{H}$ be an $L$-Lipschitz continuous and maximally monotone. Denote 
\[
S:=\zer (A+B)\neq\varnothing
\]
and assume that $(A+B)^{-1}$ is R-continuous at $0$ with modulus function $\rho$. Let $(x_k)$ be the sequence generated by the {\bf ITA} applied to the regularized problem \eqref{tseng1}. Choose $\gamma>0$ such that $\gamma^2 L^2\le 1-2\gamma\varepsilon $.
Then we have the estimate
\begin{equation}
d(x_k,S)
\leq
\delta(1+2\gamma^{-1}\epsilon^{-1})
+
\rho(\epsilon a)
\end{equation}
with the notation
\[
a:=\|\widetilde{x}\|,
\]
where $\widetilde{x}$ stands for the minimal norm element of $S$. 
By choosing
\[
\delta=\epsilon^2,
\]
we obtain the convergence condition
\begin{equation*}
d(x_k,S)
\leq
\epsilon^2+2\gamma^{-1}\epsilon+\rho(\epsilon a)\to 0 \;\; {\rm as}\; \;\epsilon\to 0.
\end{equation*}
\end{theorem}

\begin{proof}
It follows that
$$
x_{k+1}= Tx_k+t_{k+1},\;k=0,1,\ldots.
$$
Let $z_k:=x_{k}-t_k$, and let $x^*$ be the unique fixed point of $T$ that also satisfies the inclusion $0\in \tilde{A}x^*+Bx^*$. Using Lemma~\ref{ab}, we get
\begin{align*}
\|z_{k+1}-x^*\|
=\| Tx_k-x^*\|\le \kappa\|x_k-x^*\|
\leq
\kappa(\|z_k-x^*\|+\|t_k\|)\le \kappa(\|z_k-x^*\|+\delta),
\end{align*}
where $\kappa=\sqrt{1-\gamma\epsilon}<1$. The induction procedure gives us the estimates
\begin{align*}
\|z_{k+1}-x^*\|
&\leq
\kappa^{k+1}\|z_0-x^*\|
+
\kappa\delta(1+\kappa+\cdots+\kappa^k) \\
&\leq
\kappa^{k+1}\|x_0-x^*\|
+
\delta\frac{\kappa}{1-\kappa}
\end{align*}
which consequently yield the inequalities
\[
\limsup_{k\to\infty}\|z_k-x^*\|
\leq
\delta\frac{\kappa}{1-\kappa}\le\delta\frac{1+\kappa}{1-\kappa^2} \le 2\delta \gamma^{-1}\epsilon^{-1}.
\]
It is {obvious} to observe that
\[
x^*\in (A+B)^{-1}(-\epsilon x^*).
\]
Since $(A+B)^{-1}$ is R-continuous at $0$, it follows that
\[
d(x^*,S)
\leq
\rho(\epsilon\|x^*\|)\le\rho(\epsilon a).
\]
Note furthermore that
\begin{align*}
d(x_k,S)
\leq
\|x_k-z_k\|
+
\|z_k-x^*\|
+
d(x^*,S),
\end{align*}
which implies in turn that
$$
\limsup_{k\to\infty} d(x_k,S)\leq
\delta(1+
2\gamma^{-1}\epsilon^{-1})
+
\rho(\epsilon a).
$$
By choosing $\delta=\epsilon^{2}$, we complete the proof of the theorem.
\end{proof}

\end{document}
